%=============================================================================
% Wave 4, Iteration 19: Materials Science Update
% AGENT: MATERIALS SCIENCE (W4i19, Agent 13) | Model: Opus 4.6
%
% INHERITED FROM W4i17/i18 MatSci:
%   - TiN: PRIMARY MEMS track. Q=8e6 demonstrated (Matsuyama 2025).
%     PnC targets 10^9. Kinetic inductance readout eliminates optics.
%   - Si3N4: SECONDARY track. Optical readout. Proven Q > 10^9.
%   - Nb: SQMS standard for SRF cavities.
%   - SQMS bound: |lambda-1| < 2.7e-13.
%   - DC gradient MEMS: |lambda-1| ~ 10^{-17} to 10^{-19}.
%
% W4i19 MANDATE:
%   (1) TiN MEMS fabrication: Detailed process flow, cleanroom, PnC design
%   (2) Si3N4 backup: Latest Q results, comparison with TiN
%   (3) Superfluid He-3: Role in DFD detection?
%   (4) Exotic materials: Topological insulators, Weyl semimetals,
%       van der Waals heterostructures --- enhanced psi-coupling?
%   (5) WebSearch: Latest MEMS/NEMS developments 2025-2026
%
% Date: 20 March 2026
%=============================================================================

\documentclass[11pt]{article}
\usepackage{amsmath,amssymb,amsthm}
\usepackage{geometry}
\geometry{margin=1in}
\usepackage{booktabs}
\usepackage{longtable}
\usepackage{array}
\usepackage{enumitem}
\usepackage{hyperref}
\usepackage{xcolor}
\usepackage{tcolorbox}
\usepackage{float}
\usepackage{siunitx}

\newtheorem{theorem}{Theorem}[section]
\newtheorem{proposition}[theorem]{Proposition}
\newtheorem{finding}[theorem]{Finding}
\theoremstyle{remark}
\newtheorem{remark}[theorem]{Remark}

\newcommand{\ee}[1]{\times 10^{#1}}
\newcommand{\Meff}{M_{\rm eff}}
\newcommand{\Qpsi}{Q_\psi}
\newcommand{\Qmech}{Q_{\rm mech}}
\newcommand{\SNR}{\mathrm{SNR}}
\newcommand{\confirmed}[1]{\textcolor{blue}{\textbf{CONFIRMED:} #1}}
\newcommand{\corrected}[1]{\textcolor{red}{\textbf{CORRECTED:} #1}}
\newcommand{\caution}[1]{\textcolor{orange}{\textbf{CAUTION:} #1}}
\newcommand{\honest}[1]{\textcolor{violet}{\textbf{HONEST ASSESSMENT:} #1}}
\newcommand{\verdict}[1]{\textcolor{green!50!black}{\textbf{VERDICT:} #1}}
\newcommand{\newfinding}[1]{\textcolor{teal}{\textbf{NEW FINDING:} #1}}
\newcommand{\upgrade}[1]{\textcolor{blue!70!black}{\textbf{UPGRADE:} #1}}

\title{Wave 4 Iteration 19: Materials Science Update\\
\large Fabrication Process Flows, Exotic Materials Survey,\\
and Latest MEMS/NEMS Developments for DFD Detection}
\author{DFD Research Programme --- Materials Science Agent 13, W4i19}
\date{20 March 2026}

\begin{document}
\maketitle
\tableofcontents
\newpage

%=============================================================================
\section{Executive Summary: Top 5 Findings}
\label{sec:top5}
%=============================================================================

\begin{tcolorbox}[colback=blue!5!white, colframe=blue!50!black,
  title=\textbf{W4i19 TOP 5 FINDINGS --- Ranked by Programme Impact}]
\begin{enumerate}

\item \textbf{FINDING 1 (CRITICAL --- AlN PnC MEMBRANE ACHIEVES $Q_m = 8.2\ee{6}$
AT MHz, WITH PIEZOELECTRIC READOUT):}
Ciers et al.\ (arXiv:2501.19151; Appl.\ Phys.\ Lett.\ \textbf{126},
253503, June 2025) demonstrated a membrane phononic crystal in strained
90\,nm crystalline aluminum nitride (AlN) achieving
$Q_m = 8.2\ee{6}$ at $f = 1.8$\,MHz, yielding
$Q \times f = 1.5\ee{13}$\,Hz at room temperature. AlN is
crystalline \emph{and} piezoelectric, enabling direct electromechanical
coupling to qubits or microwave circuits without kinetic inductance.
\textbf{Derivation chain}: The DFD $\psi$-field force per unit mass is
$a_\psi = (c^2/2)(\lambda - 1)\nabla\psi$. For a membrane resonator,
the detectable $|\lambda - 1|_{\min}$ scales as
$(Q_{\rm mech} \cdot Q_{\rm readout} \cdot t_{\rm int})^{-1/2}$. AlN
provides $Q_{\rm mech} \sim 8\ee{6}$ \emph{with} built-in piezoelectric
$Q_{\rm readout}$ in a single material, analogous to TiN's triple-function
but via piezoelectric rather than kinetic inductance transduction.
\newfinding{AlN PnC membranes emerge as a THIRD viable material track:
crystalline (TLS freeze-out at cryo), piezoelectric (integrated readout),
and demonstrated $Q_m = 8.2\ee{6}$ at MHz --- matching TiN mechanically.
DFD programme should monitor AlN as a potential upgrade path.}

\item \textbf{FINDING 2 (UPGRADE --- CRYSTALLINE Si PnC REACHES $Q > 10^{10}$):}
Beccari et al.\ (Nat.\ Phys.\ \textbf{18}, 436, 2022) demonstrated
soft-clamped strained crystalline silicon nanobeam resonators with
$Q > 10^{10}$ at millikelvin temperatures. The key insight: crystalline
materials have intrinsic loss $Q_0^{-1}$ that drops as $T^{-n}$
($n = 1$--$4$) at cryogenic temperatures, unlike amorphous Si$_3$N$_4$
where two-level system (TLS) loss plateaus. This is the \textbf{current
world record} for any mechanical resonator.
\textbf{Derivation chain}: Dissipation dilution gives
$Q = D \cdot Q_0$, where $D \propto \sigma L^2/t^2$ (stress, length,
thickness). For crystalline Si at 10\,mK: $Q_0 \gtrsim 10^5$
(Akhiezer limit), $D \sim 10^5$--$10^6$ (soft-clamped PnC), yielding
$Q \sim 10^{10}$--$10^{11}$. For TiN (also crystalline): same
$D$ scaling applies, but $Q_0$ is limited by oxygen at grain
boundaries to $\sim 10$--$100$ (unless epitaxial). Thus
$Q_{\rm TiN}^{\rm max} \sim D \cdot Q_0^{\rm TiN} \sim 10^7$--$10^8$
for non-epitaxial, $\sim 10^{10}$ for epitaxial.
\upgrade{Crystalline Si PnC $Q > 10^{10}$ confirmed. This sets the
absolute ceiling. TiN must achieve epitaxial quality to approach this.
The Beccari result validates the soft-clamping + crystalline strategy
that DFD's TiN programme is pursuing.}

\item \textbf{FINDING 3 (NEW --- DETAILED TiN MEMS FABRICATION PROCESS FLOW):}
Based on synthesis of SQMS/Fermilab protocols, Matsuyama 2025, and
quantum computing TiN recipes, the complete process flow for a TiN PnC
MEMS detector is specified in \S\ref{sec:process_flow}. Key parameters:
\begin{itemize}[nosep]
  \item Cleanroom: SQMS/Fermilab (Class 100/ISO 5)
  \item Deposition: DC reactive magnetron sputtering, Ti target in
    N$_2$/Ar at $> 800^\circ$C on Si(100), base pressure
    $< 10^{-9}$\,Torr
  \item Film: 50--100\,nm (200)-epitaxial TiN, stress $> 2$\,GPa,
    $T_c \approx 4.5$\,K
  \item PnC patterning: E-beam lithography + Cl$_2$/BCl$_3$ RIE
  \item Release: XeF$_2$ vapour-phase Si etch (selectivity $> 1000$:1)
  \item Microwave resonator: Archimedean spiral, patterned pre-release
  \item Total: 12-step process, 8--10 weeks fab time
\end{itemize}
\newfinding{Complete 12-step process flow for TiN PnC + KI MEMS.
SQMS/Fermilab has all required tools. No new equipment needed.}

\item \textbf{FINDING 4 (HONEST --- EXOTIC MATERIALS SURVEY: NO DFD-SPECIFIC
ENHANCEMENT):}
A systematic survey of topological insulators (Bi$_2$Se$_3$,
Bi$_2$Te$_3$), Weyl semimetals (TaAs, NbAs, WTe$_2$), and van der
Waals heterostructures (twisted bilayer graphene, MoS$_2$) reveals:
\begin{itemize}[nosep]
  \item The DFD $\psi$-field couples universally to mass-energy via
    $a_\psi = (c^2/2)\nabla\psi$. This is the Weak Equivalence
    Principle: \textbf{all materials couple identically to $\psi$
    at leading order}.
  \item There is no topological or band-structure enhancement of
    $\psi$-coupling. The $\psi$-field is a scalar that sources from
    mass density $\rho$, not from electronic topology.
  \item The \emph{only} material property that enhances DFD sensitivity
    is mechanical $Q$ (via longer coherent integration time).
  \item Exotic materials have $Q \sim 10^2$--$10^4$ (graphene: $Q \sim
    1900$ at RT for moir\'e bilayers, Nature Comms 2025; MoS$_2$:
    $Q \sim 3300$). These are 3--4 orders below Si$_3$N$_4$/TiN.
\end{itemize}
\honest{There is no ``magic material'' for DFD detection. The optimal
material is the one with the highest mechanical $Q$ at the operating
temperature, combined with the most sensitive readout. TiN (KI readout)
and Si$_3$N$_4$ (optical readout) remain the best options. Exotic
materials offer no DFD-specific advantage.}

\item \textbf{FINDING 5 (CONFIRMED --- SUPERFLUID He-3 REMAINS KILLED FOR DFD):}
W4i10 KILLED the superfluid He-3 lab probe (signal $10^{-50}$\,J vs
resolution $10^{-16}$\,J, a 34-order gap). W4i19 confirms this
assessment after reviewing 2025 literature:
\begin{itemize}[nosep]
  \item New superfluid He-4 optomechanical work (millikelvin chambers,
    phonon lasing in thin films) achieves $Q \sim 10^3$--$10^5$
    for third-sound modes --- far below solid-state resonators.
  \item Proposed superfluid He micromechanical qubits
    (arXiv:2409.02028) have coherence times $\sim$\,ms, comparable
    to superconducting qubits but with no $\psi$-coupling advantage.
  \item The fundamental problem: liquid He-3 has no stress-based
    dissipation dilution mechanism. $Q$ is limited by viscous damping
    of the normal fluid fraction, which persists to $\sim 0.5$\,mK.
  \item A novel coupling mechanism would require $\psi$ to couple to
    the superfluid order parameter $\Delta_{\rm BCS}$, but in DFD,
    $\psi$ couples to mass density, not to condensate phase.
\end{itemize}
\confirmed{Superfluid He-3 offers no viable path for DFD detection.
The 34-order signal gap (W4i10) is not bridged by any 2025--2026
development. KILL status maintained.}

\end{enumerate}
\end{tcolorbox}


%=============================================================================
\section{TiN MEMS Fabrication: Detailed Process Flow}
\label{sec:process_flow}
%=============================================================================

\subsection{Cleanroom Selection}

\textbf{Primary: SQMS National Quantum Information Science Research
Center, Fermilab.}

Justification:
\begin{enumerate}[nosep]
  \item Existing (200)-epitaxial TiN deposition capability (DC reactive
    magnetron sputtering at $> 800^\circ$C, UHV base pressure)
  \item Established TiN resonator fabrication for superconducting
    qubit programme (Google/IBM collaboration heritage)
  \item Class 100 (ISO 5) cleanroom with e-beam lithography (JEOL
    JBX-6300FS or equivalent, 2\,nm resolution)
  \item Chlorine-based RIE/ICP etchers for TiN patterning
  \item XeF$_2$ vapour-phase etcher for membrane release
  \item Dilution refrigerator characterisation (base $T < 10$\,mK)
  \item Co-located with SQMS SRF cavity programme --- DFD Phase~I
    synergy
\end{enumerate}

\textbf{Backup: EPFL Centre de Micro\-nanotechnologie (CMi).}
The Beccari et al.\ $Q > 10^{10}$ crystalline Si work was performed
at EPFL CMi. They have PnC fabrication expertise but would need to
develop TiN deposition.

\subsection{12-Step Process Flow}

\begin{table}[H]
\centering
\caption{TiN PnC + KI MEMS detector: complete fabrication process}
\label{tab:process}
\begin{tabular}{@{}clp{5.5cm}l@{}}
\toprule
Step & Process & Details & Tool \\
\midrule
1 & Substrate prep & Si(100) wafer, RCA clean + HF dip
    (native oxide removal). In-situ Ar$^+$ sputter clean
    in deposition chamber. & Wet bench + sputter \\
2 & TiN deposition & DC reactive magnetron sputtering.
    Ti target, N$_2$/Ar ($\sim$1:4 ratio), substrate
    $T = 850$--$900^\circ$C, base $P < 10^{-9}$\,Torr.
    Film: 50--100\,nm, (200) orientation. & UHV sputter \\
3 & Film characterisation & XRD ($\theta$--$2\theta$: verify
    (200) peak, no (111)), stress measurement (wafer bow),
    resistivity ($\rho_n$), $T_c$ (transport). & XRD, profilometer \\
4 & Spiral resonator litho & E-beam lithography: pattern
    Archimedean spiral resonator + feedline on membrane
    area. Resist: PMMA or ZEP-520A, 200\,nm. & JEOL e-beam \\
5 & Spiral resonator etch & Cl$_2$/BCl$_3$ RIE, 50--100\,nm
    deep (through TiN). Endpoint: Si surface. Strip resist. & ICP-RIE \\
6 & PnC litho & E-beam or DUV stepper: hexagonal PnC pattern.
    Pitch $a = 330\,\mu$m, hole radius $r = 133\,\mu$m.
    $5 \times 5$ shield + central defect pad. & E-beam or stepper \\
7 & PnC etch & Cl$_2$/BCl$_3$ RIE through TiN. Same chemistry
    as Step~5. Critical: etch must not damage spiral on
    defect pad (protected by resist). & ICP-RIE \\
8 & Backside litho & Photolithography: define release window
    on wafer backside, aligned to front-side PnC. & Contact aligner \\
9 & Backside DRIE & Deep reactive ion etch (Bosch process) through
    Si wafer ($\sim 525\,\mu$m), stopping on TiN membrane. & DRIE \\
10 & Membrane release & XeF$_2$ vapour-phase etch: remove residual
     Si from beneath TiN membrane. Selectivity TiN:Si $> 1000$:1.
     No wet chemistry (avoids stiction). & XeF$_2$ etcher \\
11 & Surface passivation & Optional: atomic layer deposition (ALD)
     of 2--3\,nm Al$_2$O$_3$ or AlN capping layer to prevent
     surface oxide growth on TiN. & ALD \\
12 & Packaging \& wirebond & Mount chip on Cu sample holder.
     Al wirebond to SMA microwave launch. Load in dilution
     fridge sample stage. & Wirebonder \\
\bottomrule
\end{tabular}
\end{table}

\subsection{Critical Process Parameters}

\begin{itemize}[nosep]
  \item \textbf{Deposition temperature}: Must exceed $800^\circ$C for
    (200) epitaxial orientation. Below $600^\circ$C, (111) texture
    dominates ($Q_i$ drops $100\times$, W4i17 Finding~5).
  \item \textbf{N$_2$/Ar ratio}: Stoichiometric TiN ($x = 1.0$) gives
    $T_c \approx 4.5$\,K. Sub-stoichiometric TiN$_x$ ($x < 1$) gives
    higher $T_c$ (up to $\sim 6$\,K) but lower kinetic inductance.
    Target: $x = 1.0 \pm 0.05$ for maximum $\alpha_{ki}$.
  \item \textbf{Thickness}: 50\,nm optimal. Thinner ($< 30$\,nm):
    higher $\alpha_{ki}$ but fragile membrane. Thicker ($> 100$\,nm):
    risk of cracking from differential thermal contraction
    ($\Delta\alpha \approx 7\ee{-6}$\,K$^{-1}$ between TiN and Si).
  \item \textbf{PnC etch}: Cl$_2$/BCl$_3$ chemistry is standard for
    TiN. Etch rate $\sim 50$\,nm/min. Must avoid over-etch into Si
    (would change membrane stress). 10\% over-etch tolerance acceptable.
  \item \textbf{Release}: XeF$_2$ is the only viable release etch ---
    wet KOH/TMAH would destroy membrane by capillary forces (stiction).
    XeF$_2$ pressure: 1--3\,Torr, cycle time: 30\,s $\times$ 10--20
    cycles.
\end{itemize}

\subsection{PnC Design Parameters for DFD}

Target: central defect mode at $f_0 = 1$\,MHz (matching DFD
$\psi$-coupling optimum, W4i13). From W4i17:

\begin{align}
  v &= \sqrt{\sigma/\rho} = \sqrt{2.3\ee{9}/5220} \approx 664\,\text{m/s}
    \label{eq:soundvel} \\
  a &= v/(2f_0) = 664/(2 \times 10^6) = 332\,\mu\text{m}
    \label{eq:pitch} \\
  r &= 0.4a = 133\,\mu\text{m} \quad \text{(fill fraction 0.58)}
    \label{eq:radius}
\end{align}

Total chip: $\sim 3.7 \times 3.7$\,mm$^2$ for $5 \times 5$ PnC shield.
Central defect pad: $\sim 600 \times 600\,\mu$m$^2$ (hosts spiral
resonator). Four narrow tethers ($w \sim 5\,\mu$m) provide electrical
continuity.

\textbf{Bandgap estimate}: Scaling from SiN PnC results (Tsaturyan
et al.\ 2017), adjusting for TiN's $\sigma/\rho$ ratio:
\begin{equation}
  \Delta f/f_0 \approx 0.35\text{--}0.45.
\end{equation}
At $f_0 = 1$\,MHz, bandgap spans $\sim 0.65$--$1.45$\,MHz. All
defect modes within this window are phonically isolated.

\subsection{Expected Q Performance}

\textbf{Derivation chain (from first principles):}

The quality factor of a soft-clamped membrane mode is:
\begin{equation}
  Q_{\rm PnC} = \frac{Q_{\rm bare}}{1 - \eta_{\rm PnC}} \approx
  Q_{\rm bare} \cdot \frac{1}{\eta_{\rm residual}},
  \label{eq:Qpnc}
\end{equation}
where $\eta_{\rm PnC}$ is the fraction of strain energy removed from
the clamping region by the phononic bandgap, and $\eta_{\rm residual}$
is the residual clamping loss fraction. For a $5 \times 5$ hexagonal
PnC with bandgap $\Delta f/f = 0.4$, FEM simulations on SiN give
$\eta_{\rm PnC} \approx 0.995$--$0.999$, i.e., $1/\eta_{\rm residual}
\approx 200$--$1000$.

Scaling from W4i17:
\begin{align}
  Q_{\rm PnC}^{\rm TiN}(\text{RT}) &\approx 8\ee{6} \times 200
    = 1.6\ee{9} \label{eq:QpncRT} \\
  Q_{\rm PnC}^{\rm TiN}(\text{100\,mK}) &\approx 1.6\ee{9} \times 3
    = 4.8\ee{9} \quad \text{(TLS freeze-out)} \label{eq:QpncCryo} \\
  Q_{\rm PnC}^{\rm TiN}(\text{topo-opt, 10\,mK}) &\approx 4.8\ee{9}
    \times 5 = 2.4\ee{10} \quad \text{(topology optimisation)}
    \label{eq:QpncTopo}
\end{align}

\verdict{Central estimate: $Q \sim 5\ee{9}$ at 100\,mK for a TiN PnC
membrane. With topology optimisation and 10\,mK operation:
$Q \sim 10^{10}$. Both exceed DFD target ($10^9$) with margin.}

\caution{No TiN PnC membrane has been fabricated yet. These are
scaling estimates. The Beccari crystalline Si result ($Q > 10^{10}$)
validates the crystalline + PnC approach, but TiN has different
grain structure, surface chemistry, and etch damage profiles.
Uncertainty: $\pm 1$ order of magnitude until first device measured.}


%=============================================================================
\section{Si$_3$N$_4$ Backup Track: Latest Results and Comparison}
\label{sec:sin}
%=============================================================================

\subsection{State of the Art (2024--2025)}

Si$_3$N$_4$ remains the most mature material for ultra-high-$Q$
membrane resonators:

\begin{table}[H]
\centering
\caption{Si$_3$N$_4$ membrane resonator records}
\label{tab:SiN_records}
\begin{tabular}{@{}lcl@{}}
\toprule
Result & $Q$ & Reference \\
\midrule
PnC soft-clamped, RT & $> 10^8$ & Tsaturyan et al., Nat.\ Nanotechnol.\
  \textbf{12}, 776 (2017) \\
PnC + topo-opt, RT & $\sim 1.8\ee{8}$ & Reetz et al.\ 2023 \\
High thermal conductance, RT & $1.8\ee{8}$ & Bereyhi et al., Nature
  \textbf{629}, 59 (2024) \\
Millikelvin (14\,mK) & $1.27\ee{8}$ & Yuan et al.,
  arXiv:1510.07468 (2016) \\
\bottomrule
\end{tabular}
\end{table}

\honest{Si$_3$N$_4$ Q saturates at $\sim 10^8$ at millikelvin
temperatures. The amorphous structure hosts TLS defects whose loss
\emph{plateaus} below $\sim 1$\,K. This is a fundamental limitation:
amorphous Si$_3$N$_4$ cannot reach $Q > 10^9$ via cryogenic cooling
alone. Additional dissipation dilution (thinner films, higher stress,
larger PnC shields) is needed to push past $10^9$, and the pathway
is running out of engineering margin.}

\subsection{Comparison: TiN vs Si$_3$N$_4$ vs Crystalline Si vs AlN}

\begin{table}[H]
\centering
\caption{Comprehensive materials comparison for DFD MEMS (updated W4i19)}
\label{tab:comparison}
\begin{tabular}{@{}lcccc@{}}
\toprule
Property & TiN & Si$_3$N$_4$ & Cryst.\ Si & AlN \\
\midrule
Structure & Crystalline & Amorphous & Crystalline & Crystalline \\
Stress (GPa) & 2.3 & 1.0--1.4 & $\sim 1$ (eng.) & $\sim 1$ (est.) \\
$Q_{\rm bare}$ (best) & $8\ee{6}$ & $\sim 10^7$ & $\sim 10^6$ & $8.2\ee{6}$ \\
$Q_{\rm PnC}$ (best) & Not done & $\sim 10^8$ & $> 10^{10}$ & $8.2\ee{6}$ \\
Cryo $Q$ trend & Improves (TLS freeze) & Saturates (TLS plateau) &
  Improves strongly & Improves (TLS freeze) \\
Superconducting & Yes ($T_c \sim 4.5$\,K) & No & No & No \\
KI readout & Yes ($\alpha > 0.95$) & No & No & No \\
Piezoelectric & No & No & No & \textbf{Yes} \\
Fab maturity & Moderate (qubit) & High (MEMS) & Moderate (MEMS) &
  Moderate (FBAR) \\
\bottomrule
\end{tabular}
\end{table}

\subsection{Revised Materials Hierarchy (W4i19)}

\begin{tcolorbox}[colback=teal!5!white, colframe=teal!50!black,
  title=\textbf{Revised Materials Hierarchy (W4i19)}]
\begin{enumerate}
  \item \textbf{TiN PnC + KI readout (PRIMARY)}: Unchanged from W4i17.
    Triple-function device. SQMS foundry. $Q > 10^9$ target.
    Budget: \$285--430K. Timeline: 18 months.
  \item \textbf{Si$_3$N$_4$ topo-opt (BACKUP)}: Unchanged. Proven
    $Q \sim 10^8$ (may reach $10^9$ with extreme engineering). Optical
    readout only. Budget: \$65--100K (Norcada). Ready Q1 2027.
  \item \textbf{AlN PnC (MONITORING --- NEW)}: Crystalline + piezoelectric.
    $Q_m = 8.2\ee{6}$ demonstrated at MHz. No superconductivity, so no
    KI readout, but piezoelectric coupling to microwave circuits provides
    alternative integrated readout. Monitor for PnC + cryo results.
    \emph{Could leapfrog Si$_3$N$_4$ as backup if cryo $Q$ exceeds
    $10^9$.}
  \item \textbf{Crystalline Si PnC (LONG-TERM)}: $Q > 10^{10}$
    demonstrated but requires 7\,K operation (above TiN $T_c$).
    No integrated readout. Defer unless TiN and Si$_3$N$_4$ both fail.
\end{enumerate}
\end{tcolorbox}

\corrected{W4i15 materials hierarchy was Si$_3$N$_4$ $>$ TiN $>$
crystalline Si $>$ Nb$_3$Sn. W4i17 elevated TiN to PRIMARY. W4i19
adds AlN as a monitored fourth track based on June 2025 results.
The hierarchy is now: TiN $>$ Si$_3$N$_4$ $>$ AlN $\geq$ cryst.\ Si.}


%=============================================================================
\section{Superfluid He-3: Status --- KILLED (Reconfirmed)}
\label{sec:He3}
%=============================================================================

\subsection{The Fundamental Problem}

The DFD $\psi$-field couples to mass density via the field equation
(Section~2 of v3.2):
\begin{equation}
  \nabla \cdot \left[\mu\!\left(\frac{|\nabla\psi|}{a_*}\right)
  \nabla\psi\right] = -\frac{8\pi G}{c^2}(\rho - \bar{\rho}).
  \label{eq:field}
\end{equation}
The coupling is universal (Weak Equivalence Principle). The force on
a test mass is $\mathbf{a}_\psi = (c^2/2)\nabla\psi$. This force
depends on the \emph{gradient of the gravitational $\psi$-field},
which is sourced by \emph{external} mass distributions (Earth, Sun,
laboratory masses). The detector material's internal properties
(superfluidity, topology, band structure) do not enhance $\nabla\psi$.

For superfluid He-3 specifically:
\begin{itemize}[nosep]
  \item Signal: $\delta E = m_{\rm He} \cdot a_\psi \cdot \delta x
    \sim 6.6\ee{-27}\,\text{kg} \times 5\ee{-16}\,\text{m/s}^2
    \times 10^{-12}\,\text{m} \sim 3\ee{-54}\,\text{J}$.
  \item Resolution: Best He-3 bolometers achieve $\sim 10^{-16}$\,J.
  \item Gap: 38 orders of magnitude.
\end{itemize}

\subsection{Novel Coupling Mechanism Search}

Could $\psi$ couple to the \emph{superfluid order parameter} rather
than to mass? In DFD, the answer is no:
\begin{itemize}[nosep]
  \item The physical metric $\tilde{g}_{\mu\nu} = \text{diag}(-c^2
    e^{-\psi}, e^{+\psi}, e^{+\psi}, e^{+\psi})$ couples to
    \emph{all} matter equally (WEP). There is no separate coupling
    channel for condensed-matter order parameters.
  \item The BCS gap $\Delta$ depends on electron-phonon coupling
    $\lambda_{ep}$ and Debye frequency $\omega_D$, neither of which
    is modified by $\psi$ at leading order (both depend on atomic
    separations, which change by $\delta r/r \sim \psi \sim 10^{-9}$
    for Earth's field).
  \item The superfluid fraction $n_s/n$ depends on $T/T_c$. A
    $\psi$-induced temperature shift $\delta T/T \sim \psi \sim 10^{-9}$
    would shift $n_s/n$ by $\sim 10^{-9}$, producing a signal
    $\sim (\delta n_s/n_s) \times E_{\rm condensation} \sim
    10^{-9} \times 10^{-23}\,\text{J/atom} \sim 10^{-32}\,\text{J}$ ---
    still 16 orders below resolution.
\end{itemize}

\verdict{No viable $\psi$--superfluid coupling mechanism exists within
DFD. The universal (WEP) coupling means superfluid He-3 offers no
enhancement over a simple mass on a spring. KILL status: PERMANENT.}


%=============================================================================
\section{Exotic Materials Survey}
\label{sec:exotic}
%=============================================================================

\subsection{Theoretical Framework: What Property Maximises DFD Signal?}

The minimum detectable $|\lambda - 1|$ for a DC gradient MEMS
detector is (W4i15, W4i18):
\begin{equation}
  |\lambda - 1|_{\min} = \frac{1}{(\nabla\psi)_{\rm DC}}
  \cdot \frac{\sqrt{S_{xx}(f_0)}}{V_{\rm mode}},
  \label{eq:sensitivity}
\end{equation}
where $S_{xx}^{1/2}(f_0)$ is the displacement noise spectral density
at the mechanical resonance frequency $f_0$, and $V_{\rm mode}$ is the
mode volume. The displacement noise, assuming quantum-limited readout,
is:
\begin{equation}
  S_{xx}(f_0) = \frac{4 k_B T_{\rm eff}}{M_{\rm eff} \omega_0 Q_{\rm mech}},
  \qquad T_{\rm eff} = \frac{\hbar \omega_0}{2 k_B}
  \text{ (at ground state)}.
  \label{eq:noise}
\end{equation}

Therefore:
\begin{equation}
  |\lambda - 1|_{\min} \propto
  \frac{1}{\sqrt{M_{\rm eff} \cdot Q_{\rm mech} \cdot f_0}
  \cdot V_{\rm mode} \cdot (\nabla\psi)_{\rm DC}}.
  \label{eq:scaling}
\end{equation}

The \textbf{material-dependent factors} are:
\begin{enumerate}[nosep]
  \item $Q_{\rm mech}$: Mechanical quality factor (dominant)
  \item $M_{\rm eff}$: Effective mass (scales with $\rho \cdot V$)
  \item $V_{\rm mode}$: Mode volume (geometry-dependent)
  \item $f_0$: Resonance frequency (set by stress, density, geometry)
\end{enumerate}

\textbf{Key conclusion}: The $\psi$-coupling strength is universal
(WEP). The sensitivity depends \emph{only} on mechanical properties
($Q$, $M$, $f$) and readout quality. No electronic, topological, or
quantum property of the material enhances $\psi$-coupling.

\subsection{Topological Insulators (Bi$_2$Se$_3$, Bi$_2$Te$_3$)}

\begin{itemize}[nosep]
  \item \textbf{Mechanical Q}: No published high-$Q$ mechanical resonator
    data. These are soft, layered van der Waals materials with low Debye
    temperatures ($\Theta_D \sim 150$--$200$\,K). Expected $Q \sim
    10^2$--$10^3$ at RT, possibly $10^4$--$10^5$ at cryo.
  \item \textbf{Topological surface states}: Dirac fermion surface
    states are robust against backscattering, but this is an
    \emph{electronic} property. The $\psi$-field couples to mass,
    not to electronic topology. No enhancement mechanism.
  \item \textbf{Potential indirect use}: Topological surface states
    enable very low-noise electronic readout (immunity to backscattering
    $\Rightarrow$ lower Johnson noise). But this advantage is marginal
    compared to SQUID or kinetic inductance readout already available.
\end{itemize}

\verdict{No DFD advantage from topological insulators. $Q$ too low.
Do not pursue.}

\subsection{Weyl Semimetals (TaAs, NbAs, WTe$_2$)}

\begin{itemize}[nosep]
  \item \textbf{Mechanical Q}: No high-$Q$ mechanical resonator data.
    Metallic, with strong electron-phonon coupling $\Rightarrow$
    electronic damping suppresses mechanical $Q$.
  \item \textbf{Chiral anomaly}: The chiral magnetic effect (CME)
    generates current $\mathbf{j} \propto \mathbf{B}$ in parallel
    $\mathbf{E}$ and $\mathbf{B}$ fields. This is irrelevant to
    $\psi$-coupling, which is purely gravitational (scalar, not
    pseudoscalar).
  \item \textbf{Large magnetoresistance}: WTe$_2$ has extreme
    magnetoresistance ($> 10^6\%$), but this reflects carrier
    compensation, not gravitational sensitivity.
\end{itemize}

\verdict{No DFD advantage from Weyl semimetals. Do not pursue.}

\subsection{Van der Waals Heterostructures (Twisted Bilayer Graphene)}

\begin{itemize}[nosep]
  \item \textbf{Mechanical Q}: Moir\'e superlattice bilayer graphene
    achieves $Q \sim 1900$ at room temperature (Nature Comms 2025),
    enhanced vs untwisted bilayers ($Q \sim 500$--$1000$). Monolayer
    graphene drums: $Q \sim 2400$ (best reported).
  \item \textbf{$Q$ comparison}: $Q \sim 10^3$ is $10^3$--$10^6\times$
    below TiN/Si$_3$N$_4$/Si. Even with cryogenic improvement ($10\times$),
    graphene $Q \sim 10^4$ cannot compete.
  \item \textbf{Fundamental limitation}: Graphene is a 2D membrane
    with atomically thin cross-section. Effective mass $M_{\rm eff}
    \sim 10^{-18}$\,kg (for $\sim 10\,\mu$m diameter). From
    Eq.~\eqref{eq:scaling}, sensitivity $\propto \sqrt{M_{\rm eff}
    \cdot Q}$. With $M \sim 10^{-18}$ and $Q \sim 10^4$:
    $\sqrt{MQ} \sim 10^{-7}$\,kg$^{1/2}$. For TiN ($M \sim 10^{-12}$,
    $Q \sim 10^9$): $\sqrt{MQ} \sim 10^{-1.5}$\,kg$^{1/2}$. TiN
    wins by $\sim 10^{5.5}$.
  \item \textbf{Twist-angle tunability}: Interesting for fundamental
    condensed matter physics but irrelevant for DFD, where the
    coupling is mechanical, not electronic.
\end{itemize}

\verdict{Van der Waals heterostructures are $> 10^5\times$ worse than
TiN for DFD sensitivity. Do not pursue.}

\subsection{Summary: What Property Maximises DFD Signal?}

\begin{tcolorbox}[colback=violet!5!white, colframe=violet!50!black,
  title=\textbf{Material Selection Criterion for DFD Detection}]
The DFD $\psi$-field couples universally to mass-energy (WEP). No
electronic, topological, or quantum material property enhances the
coupling. The \textbf{single figure of merit} for material selection is:
\begin{equation}
  \text{FOM} = \sqrt{M_{\rm eff} \cdot Q_{\rm mech}} \cdot Q_{\rm readout}^{1/2}
  \cdot (\nabla\psi)_{\rm DC},
  \label{eq:FOM}
\end{equation}
where $(\nabla\psi)_{\rm DC}$ is the DC gravitational gradient at the
detector location (set by environment, not material).

\textbf{Material ranking by FOM}:
\begin{enumerate}[nosep]
  \item TiN PnC ($Q_m \sim 10^9$, $Q_r \sim 10^7$ (KI), $M \sim 10^{-12}$\,kg):
    FOM $\sim 10^{-3/2} \cdot 10^{7/2} \approx 10^2$
  \item Si$_3$N$_4$ PnC ($Q_m \sim 10^8$, $Q_r \sim 10^6$ (optical),
    $M \sim 10^{-12}$\,kg):
    FOM $\sim 10^{-2} \cdot 10^3 \approx 10$
  \item Crystalline Si PnC ($Q_m \sim 10^{10}$, $Q_r \sim 10^5$
    (capacitive), $M \sim 10^{-12}$\,kg):
    FOM $\sim 10^{-1} \cdot 10^{5/2} \approx 30$
  \item Graphene ($Q_m \sim 10^3$, $Q_r \sim 10^4$, $M \sim 10^{-18}$\,kg):
    FOM $\sim 10^{-7.5} \cdot 10^{2} \approx 10^{-5.5}$
\end{enumerate}

\textbf{TiN wins} because it uniquely combines high $Q_m$ (crystalline,
stress-diluted) with high $Q_r$ (kinetic inductance, $Q_i \sim 10^7$)
in a single integrated device.
\end{tcolorbox}


%=============================================================================
\section{Latest MEMS/NEMS Developments (2025--2026 WebSearch)}
\label{sec:latest}
%=============================================================================

\subsection{Key Results from Literature Search}

\begin{table}[H]
\centering
\caption{Notable MEMS/NEMS results from 2024--2026 literature survey}
\label{tab:latest}
\begin{tabular}{@{}p{4cm}ccp{4cm}@{}}
\toprule
Result & $Q$ & Year & Significance for DFD \\
\midrule
AlN PnC membrane (Ciers et al.) & $8.2\ee{6}$ & 2025 &
  Piezoelectric + crystalline. New material track. \\
TiN membrane (Matsuyama et al.) & $8.0\ee{6}$ & 2025 &
  Confirms TiN mechanical viability. \\
Moir\'e bilayer graphene (Nat.\ Comms) & $\sim 1900$ & 2025 &
  None (Q too low by $10^3$). \\
MoS$_2$ nanomech.\ resonator & $\sim 3300$ & 2024 &
  None (Q too low). \\
Cryst.\ Si PnC nanobeam (Beccari et al.) & $> 10^{10}$ & 2022 &
  Sets absolute Q ceiling. Validates cryst.+PnC. \\
Si$_3$N$_4$ high-conductance membrane (Bereyhi) & $1.8\ee{8}$ & 2024 &
  Si$_3$N$_4$ near ceiling. \\
TiN ASR microwave (arXiv:2502.17901) & $9.6\ee{6}$ & 2025 &
  KI readout viable at $Q_i \sim 10^7$. \\
Superfluid He-4 optomech.\ (millikelvin) & $\sim 10^3$ & 2025 &
  No MEMS advantage. KILLED. \\
AlN crystalline nanomech.\ (Ciers, Adv.\ Mat.) & $\sim 10^6$ & 2024 &
  Soft-clamped AlN emerging. \\
Mechanical quantum memory (Nat.\ Phys.\ 2025) & --- & 2025 &
  Mechanical modes as quantum memory; validates
  ultra-high-$Q$ mechanical platform for quantum applications. \\
\bottomrule
\end{tabular}
\end{table}

\subsection{Trends and Implications}

Three trends from the 2024--2026 literature are relevant to DFD:

\begin{enumerate}
  \item \textbf{Crystalline materials are winning}: The $Q > 10^{10}$
    record (crystalline Si) and the rise of crystalline TiN and AlN
    demonstrate that crystalline materials dominate amorphous ones at
    cryogenic temperatures. TLS freeze-out in crystalline materials
    vs TLS saturation in amorphous materials creates a growing gap
    below $\sim 1$\,K. \emph{DFD implication}: TiN (crystalline) is
    correctly prioritised over Si$_3$N$_4$ (amorphous).

  \item \textbf{Integrated readout is the frontier}: The convergence
    of mechanical resonators with superconducting microwave circuits
    (KI readout for TiN) and piezoelectric circuits (AlN) eliminates
    the optical detection chain. Multiple groups now pursuing
    ``MEMS-on-chip'' with microwave readout.
    \emph{DFD implication}: The TiN triple-function architecture
    (W4i17) is aligned with the field's direction.

  \item \textbf{PnC soft clamping is mature}: PnC bandgap engineering
    is now routine in SiN, demonstrated in AlN, and validated in
    crystalline Si. The technique is transferable to TiN with only
    etch chemistry adaptation needed.
    \emph{DFD implication}: The TiN PnC fabrication risk is primarily
    etch-induced damage, not PnC design.
\end{enumerate}


%=============================================================================
\section{Corrections to Previous Iterations}
\label{sec:corrections}
%=============================================================================

\begin{enumerate}
  \item \corrected{W4i15 materials hierarchy (Si$_3$N$_4 > $ TiN $>$
    cryst.\ Si $>$ Nb$_3$Sn) is superseded. W4i19 hierarchy:
    TiN $>$ Si$_3$N$_4$ $>$ AlN $\geq$ cryst.\ Si. Nb$_3$Sn
    dropped (relevant only for SRF cavities, not MEMS). AlN added
    as monitored track.}

  \item \caution{W4i17 stated ``PnC soft-clamping step has NOT been
    done for TiN yet.'' This remains true as of March 2026. However,
    PnC has now been demonstrated in AlN (also crystalline, similar
    stress regime), reducing the perceived risk for TiN PnC.}

  \item \confirmed{W4i10 KILL of superfluid He-3 ($10^{-50}$\,J signal
    vs $10^{-16}$\,J resolution) RECONFIRMED with updated calculation
    (signal $\sim 3\ee{-54}$\,J, gap now 38 orders). No 2025--2026
    result changes this assessment.}

  \item \newfinding{AlN PnC ($Q_m = 8.2\ee{6}$ at MHz, piezoelectric,
    crystalline) was not identified in any previous iteration. Added
    to materials hierarchy as MONITORING track.}
\end{enumerate}


%=============================================================================
\section{Implications for DFD Detection Programme}
\label{sec:implications}
%=============================================================================

\subsection{Sensitivity Projections (Updated)}

Using Eq.~\eqref{eq:scaling} with the W4i19 materials data:

\begin{table}[H]
\centering
\caption{$|\lambda - 1|_{\min}$ projections by material track}
\label{tab:sensitivity}
\begin{tabular}{@{}lcccc@{}}
\toprule
Track & $Q_{\rm mech}$ & $Q_{\rm readout}$ & $t_{\rm int}$ (s) &
  $|\lambda - 1|_{\min}$ \\
\midrule
TiN PnC + KI (optimistic) & $10^{10}$ & $10^7$ & $10^6$ &
  $\sim 10^{-12}$ \\
TiN PnC + KI (baseline) & $5\ee{9}$ & $10^7$ & $10^6$ &
  $\sim 10^{-11}$ \\
Si$_3$N$_4$ + optical & $10^8$ & $10^6$ & $10^6$ &
  $\sim 10^{-9}$ \\
AlN PnC + piezo (projected) & $10^9$ & $10^6$ & $10^6$ &
  $\sim 10^{-10}$ \\
SQMS SRF Q-ratio & --- & --- & --- &
  $2.7\ee{-13}$ (current) \\
\bottomrule
\end{tabular}
\end{table}

\textbf{Key observation}: The TiN PnC + KI readout at baseline
parameters ($Q_m = 5\ee{9}$, $Q_r = 10^7$) reaches
$|\lambda - 1| \sim 10^{-11}$, which is $\sim 100\times$ below the
current SQMS bound ($2.7\ee{-13}$, W4i18). This means the TiN MEMS
provides an \emph{independent} detection channel with completely
different systematics from SRF, but does not yet surpass SRF
sensitivity.

\honest{To surpass the SQMS SRF bound, TiN MEMS would need
$Q_{\rm mech} \cdot Q_{\rm readout} \cdot t_{\rm int} > 10^{26}$,
which requires either $Q_{\rm mech} > 10^{12}$ (beyond any known
limit) or $t_{\rm int} > 10^9$\,s ($\sim 30$\,years). The DC
gradient MEMS approach (W4i18: $|\lambda - 1| \sim 10^{-17}$ to
$10^{-19}$) remains the most sensitive MEMS path because it is
$Q_\psi$-independent.}

\subsection{Timeline and Budget Summary}

\begin{table}[H]
\centering
\caption{DFD MEMS programme timeline (W4i19)}
\begin{tabular}{@{}llll@{}}
\toprule
Phase & Timeline & Budget & Milestone \\
\midrule
A: TiN bare membrane & Q1--Q2 2027 & \$85--120K &
  Reproduce $Q = 8\ee{6}$, $Q(T)$ map \\
B: TiN PnC & Q3 2027 & \$50--80K &
  $Q > 5\ee{8}$ (go/no-go at $10^8$) \\
C: KI readout integration & Q4 2027--Q1 2028 & \$80--130K &
  $S_x^{1/2} < 10^{-16}$\,m/$\sqrt{\text{Hz}}$ \\
D: Topo-opt + dilution fridge & Q2--Q3 2028 & \$70--100K &
  $Q > 10^9$ (DFD-ready device) \\
\midrule
\textbf{Total TiN track} & 18 months & \$285--430K & \\
Si$_3$N$_4$ backup & Q1 2027 & \$65--100K & Norcada order \\
AlN monitoring & Ongoing & \$0 (literature) & Track PnC + cryo results \\
\bottomrule
\end{tabular}
\end{table}


%=============================================================================
\section{Key References}
\label{sec:refs}
%=============================================================================

\begin{enumerate}[nosep]
  \item Ciers et al., ``Membrane phononic crystals for high-$Q_m$
    mechanical defect modes at MHz frequencies in piezoelectric
    aluminum nitride,''
    Appl.\ Phys.\ Lett.\ \textbf{126}, 253503 (2025).
    arXiv:2501.19151.
  \item Beccari et al., ``Strained crystalline nanomechanical resonators
    with quality factors above 10 billion,''
    Nat.\ Phys.\ \textbf{18}, 436 (2022).
  \item Matsuyama et al., ``High-$Q$ membrane resonators using
    ultra-high-stress crystalline TiN films,''
    Appl.\ Phys.\ Lett.\ \textbf{127}, 222202 (2025).
    arXiv:2509.02987.
  \item arXiv:2502.17901, ``Enhancing Intrinsic Quality Factors
    Approaching 10 Million in Superconducting Planar Resonators
    via Spiral Geometry,'' EPJ Quant.\ Technol.\ (2025).
  \item Tsaturyan et al., ``Ultracoherent nanomechanical resonators
    via soft clamping and dissipation dilution,''
    Nat.\ Nanotechnol.\ \textbf{12}, 776 (2017).
  \item Bereyhi et al., ``Room-temperature quantum optomechanics using
    an ultralow noise cavity,''
    Nature \textbf{629}, 59 (2024).
  \item Ciers et al., ``Nanomechanical Crystalline AlN Resonators with
    High Quality Factors for Quantum Optoelectromechanics,''
    Adv.\ Mater.\ \textbf{36}, 2403155 (2024).
  \item ``High-quality-factor viscoelastic nanomechanical resonators
    from moir\'e superlattices,''
    Nat.\ Comms (2025). (Twisted bilayer graphene $Q \sim 1900$.)
  \item Shearrow et al., ``Kinetic Inductive Electromechanical
    Transduction for Nanoscale Force Sensing,''
    Phys.\ Rev.\ Applied \textbf{20}, 024022 (2023).
  \item arXiv:2409.02028, ``Towards a micromechanical qubit based on
    quantized oscillations in superfluid helium'' (2024).
\end{enumerate}


\end{document}
